\documentclass[11pt,a4paper]{article}
\usepackage[T2A]{fontenc}
\usepackage[utf8]{inputenc}
\usepackage[russian]{babel}
\usepackage{amsmath,amssymb,amsfonts}
\usepackage{geometry}
\geometry{top=20mm,bottom=20mm,left=20mm,right=20mm}
\usepackage{graphicx}
\usepackage{booktabs}
\usepackage{tabularx}
\usepackage{array}
\usepackage{xcolor}
\usepackage{hyperref}
\usepackage{listings}

% Цветовая палитра оформления
\definecolor{primary}{RGB}{20, 50, 110}
\definecolor{secondary}{RGB}{0, 102, 102}
\definecolor{accent}{RGB}{160, 40, 40}
\definecolor{codebg}{RGB}{248, 249, 250}
\definecolor{boxbg}{RGB}{240, 244, 250}
\definecolor{boxframe}{RGB}{70, 100, 160}
\definecolor{warnbg}{RGB}{255, 248, 240}
\definecolor{warnframe}{RGB}{210, 120, 30}
\definecolor{darkgreen}{RGB}{20, 120, 40}

\hypersetup{
    colorlinks=true,
    linkcolor=primary,
    citecolor=secondary,
    urlcolor=accent,
    pdftitle={Численное интегрирование методом Монте-Карло},
    pdfauthor={МГТУ им. Н.Э. Баумана, каф. ФН-4}
}

% Настройка листингов кода
\lstset{
    language=C++,
    backgroundcolor=\color{codebg},
    basicstyle=\ttfamily\footnotesize,
    keywordstyle=\color{primary}\bfseries,
    commentstyle=\color{darkgreen}\itshape,
    stringstyle=\color{accent},
    numbers=left,
    numberstyle=\tiny\color{gray},
    stepnumber=1,
    numbersep=8pt,
    tabsize=4,
    showspaces=false,
    showstringspaces=false,
    breaklines=true,
    frame=single,
    rulecolor=\color{boxframe},
    captionpos=b
}

\newcommand{\infobox}[2]{%
  \par\vspace{3mm}\noindent
  \fcolorbox{boxframe}{boxbg}{%
    \parbox{\dimexpr\textwidth-2\fboxsep-2\fboxrule\relax}{%
      \textbf{\color{primary}#1}\vspace{1.5mm}\par\small #2
    }%
  }\par\vspace{3mm}
}

\newcommand{\warnbox}[2]{%
  \par\vspace{3mm}\noindent
  \fcolorbox{warnframe}{warnbg}{%
    \parbox{\dimexpr\textwidth-2\fboxsep-2\fboxrule\relax}{%
      \textbf{\color{accent}#1}\vspace{1.5mm}\par\small #2
    }%
  }\par\vspace{3mm}
}

\title{\vspace{-10mm}\Huge\bfseries\color{primary} Численное интегрирование методом Монте-Карло и методы понижения дисперсии\\[3mm]
\Large\color{secondary} Теоретическое руководство и сравнительный анализ к лабораторной работе №6}
\author{\textbf{Кафедра «Физика» (ФН-4), МГТУ им. Н.Э. Баумана}\\
Курс «Вычислительная физика» (7 семестр)}
\date{2026}

\begin{document}

\maketitle

\begin{abstract}
\noindent Настоящее теоретическое руководство посвящено методу Монте-Карло (методу статистических испытаний) в контексте вычисления определенных интегралов. Подробно рассмотрены классификация задач, специфика и фундаментальные особенности метода, общая схема вероятностного моделирования, геометрический метод («Hit-or-Miss») и метод выборочного среднего («Sample Mean»). Доказано строгое превосходство метода выборочного среднего по дисперсии. Исчерпывающе изложен математический аппарат методов понижения дисперсии: метода противоположных (антитетических) переменных, метода контрольных переменных, расслоенной выборки и выборки по важности. 

Приведены полные аналитические вычисления интегралов и теоретических дисперсий для практических заданий работы ($\int_0^1 e^x dx$ и $\int_0^{\pi/2} \sin(x) dx$), а также для интегралов предыдущей лабораторной работы №5 (\textbf{Вариант 7}: $\int_0^{\sqrt{3}}\sqrt{3-x^2}dx$ и $\int_0^{3\pi/4}x\sin(2x/3)dx$). Представлен детальный сравнительный анализ статистического метода Монте-Карло и классических квадратур Ньютона--Котеса (прямоугольников, трапеций, Симпсона и экстраполяции Ричардсона). Даны рекомендации по численно устойчивой реализации алгоритмов на языке C++.
\end{abstract}

\tableofcontents
\newpage

\section{Теоретические основы метода Монте-Карло}

\subsection{1. Задачи, решаемые методом Монте-Карло}

Метод Монте-Карло (ММК) представляет собой численный метод решения широкого класса математических и прикладных физических задач, базирующийся на моделировании случайных величин и статистической оценке их характеристик.

Исторически метод зародился в конце 1940-х годов в Лос-Аламосской национальной лаборатории (США) благодаря работам С.~Улама, Дж.~фон Неймана и Н.~Метрополиса в рамках ядерного проекта. Все решаемые задачи классифицируются на две фундаментальные группы:

\subsubsection{А. Стохастические задачи (прямое физическое моделирование)}
В данных задачах исследуемый физический процесс носит принципиально вероятностный (квантовомеханический или термодинамический) характер. ММК выступает в качестве численного эксперимента, прямо воспроизводящего случайную микрофизику:
\begin{enumerate}
    \item \textbf{Перенос излучения и нейтронная физика:} Моделирование индивидуальных траекторий нейтронов и гамма-квантов в активных зонах ядерных реакторов, расчет коэффициента размножения $k_{eff}$, проектирование радиационной защиты, детектирование элементарных частиц (пакеты Geant4, MCNP).
    \item \textbf{Статистическая физика и физическая кинетика:} Метод Монте-Карло по схеме марковских цепей (MCMC, алгоритм Метрополиса--Гастингса) для канонического и микроканонического ансамблей; моделирование спиновых систем (модель Изинга, модель Гейзенберга), фазовых переходов, расчет уравнений состояния газов и жидкостей.
    \item \textbf{Оптика мутных сред и биофотоника:} Моделирование многократного рассеяния лазерного излучения в биологических тканях, атмосфере и океане.
    \item \textbf{Теория перколяции и полимеры:} Исследование критических явлений, проводимости неоднородных сред, конформаций макромолекул.
    \item \textbf{Теория надежности и массового обслуживания:} Оценка вероятностей редких катастрофических событий при эксплуатации сложных инженерных комплексов.
\end{enumerate}

\subsubsection{Б. Детерминированные задачи (математический анализ и математическая физика)}
В данных задачах исходная постановка не содержит элементов случайности, однако задача искусственно отображается на эквивалентную вероятностную модель:
\begin{enumerate}
    \item \textbf{Вычисление многомерных интегралов:} Интегрирование по областям произвольной формы в пространствах размерности $d \ge 4 \dots 100$. В таких размерностях классические сеточные методы неприменимы из-за комбинаторного взрыва числа узлов.
    \item \textbf{Решение систем линейных алгебраических уравнений (СЛАУ):} Метод случайных блужданий по цепочкам индексов матрицы для нахождения решений систем вида $(I - A)x = f$. Преимущество ММК заключается в возможности найти одну конкретную компоненту вектора $x_i$ без обращения всей матрицы и без нахождения остальных компонент.
    \item \textbf{Краевые задачи для дифференциальных уравнений в частных производных (ДУЧП):} Решение задачи Дирихле для уравнений эллиптического типа ($\Delta u = 0$, $\Delta u = -f$) методом «блуждания по сферам» (Walk-on-Spheres) или случайных блужданий по решетке, основанное на стохастическом представлении решений уравнениями Ито и формулой Фейнмана--Каца.
    \item \textbf{Глобальная оптимизация в многомерных пространствах:} Поиск глобальных экстремумов сложных многоэкстремальных невыпуклых функций (метод имитации отжига -- Simulated Annealing).
\end{enumerate}

---

\subsection{2. Основные особенности метода Монте-Карло}

\begin{enumerate}
    \item \textbf{Статистическая природа оценки:}
    Результат вычислений не является точным детерминированным числом. Он представляет собой реализацию случайной величины $\bar{\xi}_N$, сходящейся к истинному результату по закону больших чисел. Результат обязательно должен сопровождаться доверительным интервалом.
    
    \item \textbf{Универсальная скорость сходимости $O(N^{-1/2})$:}
    В силу Центральной предельной теоремы (ЦПТ) среднеквадратичная погрешность метода убывает обратно пропорционально квадратному корню из числа независимых испытаний:
    \begin{equation}
        \varepsilon_N \approx \frac{\sigma}{\sqrt{N}},
        \label{eq:clt_rate}
    \end{equation}
    где $\sigma$ --- генеральное стандартное отклонение. Чтобы увеличить точность на один десятичный знак (уменьшить погрешность в 10 раз), требуется увеличить объем выборки $N$ в $100$ раз.
    
    \item \textbf{Независимость скорости сходимости от размерности пространства $d$:}
    Погрешность классических сеточных квадратурных формул порядка точности $p$ в $d$-мерном пространстве составляет $O(N^{-p/d})$ (где $N = m^d$ --- общее число узлов сетки). При возрастании $d$ сходимость сеточных методов катастрофически деградирует («проклятие размерности»). Напротив, погрешность ММК сохраняет асимптотику $O(N^{-1/2})$ при любой размерности $d \in \mathbb{N}$.
    
    \item \textbf{Инвариантность к геометрической форме области интегрирования:}
    В отличие от сеточных методов, требующих конформных сеток или параметризации криволинейных границ, метод Монте-Карло требует лишь программной реализации индикаторной функции (предиката) вида \texttt{is\_inside(x)}.
    
    \item \textbf{Внутренняя апостериорная оценка погрешности «на лету»:}
    Оценка дисперсии $s^2$ и доверительного интервала накапливается в процессе того же самого статистического прогона без необходимости повторных вычислений (в отличие от правила Рунге в сеточных методах).
    
    \item \textbf{Критическая чувствительность к генераторам псевдослучайных чисел (ГПСЧ):}
    Качество результата жестко лимитировано периодом генератора, равномерностью распределения в многомерных кубах и отсутствием межвыборочных автокорреляций. Устаревший генератор \texttt{rand()} непригоден; стандартом вычислительной физики являются 64-битный вихрь Мерсенна (\texttt{std::mt19937\_64}) или алгоритмы семейства PCG/Xoshiro.
\end{enumerate}

---

\subsection{3. Общая схема метода Монте-Карло}

Общий процесс вычисления детерминированной величины $I$ методом статистических испытаний состоит из четырех канонических этапов:

\begin{figure}[h!]
\centering
\fbox{
\begin{minipage}{0.92\textwidth}
\textbf{Алгоритмическая схема ММК:}
\begin{enumerate}
    \item \textbf{Построение вероятностной модели:} Представление $I = \mathbb{E}[\xi] = \int g(x) p(x) dx$.
    \item \textbf{Моделирование случайной величины $\xi$:} Преобразование равномерных псевдослучайных чисел $U \sim \mathcal{U}(0,1)$ в значения случайного аргумента $X \sim p(x)$.
    \item \textbf{Статистический эксперимент:} Проведение $N$ независимых испытаний, получение выборки $\xi_1, \xi_2, \dots, \xi_N$.
    \item \textbf{Статистическая обработка:} Расчет среднего $\bar{\xi}_N$, выборочной дисперсии $s^2$ и доверительного интервала по ЦПТ.
\end{enumerate}
\end{minipage}
}
\end{figure}

\subsubsection{Моделирование случайных величин}
Для генерации случайных чисел с заданной плотностью $p(x)$ из базовых равномерных величин $U \sim \mathcal{U}(0,1)$ применяются фундаментальные методы:
\begin{itemize}
    \item \textbf{Метод обратной функции распределения (Inverse Transform Method):}
    Пусть $F(x) = \int_{-\infty}^x p(t) dt$ --- непрерывная строго монотонная функция распределения. Тогда случайная величина
    \begin{equation}
        X = F^{-1}(U), \quad U \sim \mathcal{U}(0,1)
    \end{equation}
    имеет плотность вероятности $p(x)$.
    
    \item \textbf{Метод исключения (метод фон Неймана, Rejection Sampling):}
    Если плотность $p(x)$ мажорируется вспомогательной плотностью $q(x)$ (для которой известен метод обратной функции) с константой $M \ge 1$: $p(x) \le M q(x)$, то генерируется пара $X \sim q(x)$ и $U \sim \mathcal{U}(0,1)$. Точка $X$ принимается, если $U \le \frac{p(X)}{M q(X)}$, и отбрасывается в противном случае.
\end{itemize}

\subsubsection{Оценка математического ожидания и доверительные интервалы}
По закону больших чисел в форме Хинчина несмещенная точечная оценка математического ожидания равна среднему арифметическому выборки:
\begin{equation}
    \bar{\xi}_N = \frac{1}{N} \sum_{i=1}^N \xi_i \xrightarrow[N \to \infty]{\text{п.н.}} I = \mathbb{E}[\xi].
\end{equation}

Несмещенная выборочная дисперсия случайной величины $\xi$:
\begin{equation}
    s^2 = \frac{1}{N - 1} \sum_{i=1}^N (\xi_i - \bar{\xi}_N)^2 = \frac{1}{N - 1} \left( \sum_{i=1}^N \xi_i^2 - \frac{1}{N} \left( \sum_{i=1}^N \xi_i \right)^2 \right).
    \label{eq:sample_var}
\end{equation}

Дисперсия выборочного среднего $\bar{\xi}_N$:
\begin{equation}
    \mathrm{Var}(\bar{\xi}_N) = \frac{\mathrm{Var}(\xi)}{N} \approx \frac{s^2}{N}, \quad S_{\bar{\xi}} = \frac{s}{\sqrt{N}}.
\end{equation}

Согласно Центральной предельной теореме Ляпунова, при $N \to \infty$ нормированное отклонение выборочного среднего распределено по стандартному нормальному закону:
\begin{equation}
    \frac{\bar{\xi}_N - I}{s / \sqrt{N}} \xrightarrow{d} \mathcal{N}(0, 1).
\end{equation}

Следовательно, с доверительной вероятностью $\gamma = 1 - \alpha$ истинное значение интеграла $I$ локализовано в границах:
\begin{equation}
    I \in \left[ \bar{\xi}_N - z_{1 - \alpha/2} \frac{s}{\sqrt{N}}, \quad \bar{\xi}_N + z_{1 - \alpha/2} \frac{s}{\sqrt{N}} \right],
\end{equation}
где $z_{1 - \alpha/2}$ --- квантиль стандартного нормального распределения. В частности:
\begin{itemize}
    \item Для доверительной вероятности $\gamma = 0.95$ ($\alpha = 0.05$): $z_{0.975} \approx 1.960$.
    \item Для правила «трех сигм» ($\gamma = 0.9973$): $z \approx 3.000$, предельная абсолютная погрешность $\varepsilon_{99.7\%} = 3 \frac{s}{\sqrt{N}}$.
\end{itemize}

---

\subsection{4. Вычисление определенных интегралов и методы понижения дисперсии}

Рассматривается задача вычисления одномерного определенного интеграла:
\begin{equation}
    I = \int_{a}^{b} f(x)\,dx.
    \label{eq:integral_base}
\end{equation}

\subsubsection{А. Простейшие методы}

\paragraph{1. Геометрический метод («Hit-or-Miss», метод попаданий и промахов)}
Пусть функция неотрицательна и ограничена: $0 \le f(x) \le M$ на отрезке $[a,b]$. Поместим криволинейную трапецию под графиком $f(x)$ в прямоугольник $D = [a,b] \times [0, M]$ площадью $S_D = M(b - a)$.

Генерируются независимые случайные векторы $(X_i, Y_i)$, равномерно распределенные в прямоугольнике $D$:
\begin{equation}
    X_i = a + (b - a) U_{1,i}, \quad Y_i = M \cdot U_{2,i}, \quad U_{1,i}, U_{2,i} \sim \mathcal{U}(0,1).
\end{equation}

Вводится индикатор попадания точки под кривую:
\begin{equation}
    \chi_i = \mathbb{I}(Y_i \le f(X_i)) = \begin{cases} 1, & Y_i \le f(X_i), \\ 0, & Y_i > f(X_i). \end{cases}
\end{equation}

Вероятность попадания точки под кривую равна отношению площадей:
\begin{equation}
    p = \mathbb{P}(Y \le f(X)) = \frac{I}{S_D} = \frac{I}{M(b - a)}.
\end{equation}

Случайная величина $\chi_i$ подчиняется распределению Бернулли с параметром $p$. Ее математическое ожидание $\mathbb{E}[\chi] = p$, а дисперсия $\mathrm{Var}(\chi) = p(1 - p)$.

Оценка интеграла по выборке объема $N$:
\begin{equation}
    I_{geom} = S_D \cdot \frac{N_{hit}}{N} = \frac{M(b - a)}{N} \sum_{i=1}^N \chi_i.
\end{equation}

Дисперсия оценки геометрического метода:
\begin{equation}
    \sigma^2_{geom} = \mathrm{Var}(I_{geom}) = S_D^2 \frac{p(1 - p)}{N} = \frac{I \big( M(b - a) - I \big)}{N}.
    \label{eq:var_geom}
\end{equation}

\paragraph{2. Метод выборочного среднего («Sample Mean»)}
Интеграл \eqref{eq:integral_base} тождественно переписывается через плотность непрерывного равномерного распределения $p(x) = \frac{1}{b - a}$ на $[a,b]$:
\begin{equation}
    I = (b - a) \int_{a}^{b} f(x) \frac{1}{b - a} dx = (b - a) \mathbb{E}[f(X)], \quad X \sim \mathcal{U}(a,b).
\end{equation}

Оценка метода выборочного среднего по $N$ точкам $X_i = a + (b - a)U_i$:
\begin{equation}
    I_{mean} = \frac{b - a}{N} \sum_{i=1}^N f(X_i).
    \label{eq:sample_mean_est}
\end{equation}

Дисперсия этой оценки:
\begin{equation}
    \mathrm{Var}(I_{mean}) = \frac{(b - a)^2}{N} \mathrm{Var}(f(X)) = \frac{1}{N} \left[ (b - a) \int_{a}^{b} f^2(x) dx - I^2 \right].
    \label{eq:var_mean}
\end{equation}

\infobox{Теорема о превосходстве метода выборочного среднего над геометрическим}{
Для любой неотрицательной ограниченной функции $0 \le f(x) \le M$ на отрезке $[a,b]$ дисперсия метода выборочного среднего всегда строго не превосходит дисперсии геометрического метода:
$$\sigma^2_{mean} \le \sigma^2_{geom}.$$
\textit{Доказательство:}
Вычтем дисперсию \eqref{eq:var_mean} из дисперсии \eqref{eq:var_geom}:
\begin{align*}
N \cdot \big(\sigma^2_{geom} - \sigma^2_{mean}\big) &= \big[ I \cdot M(b - a) - I^2 \big] - \left[ (b - a)\int_{a}^{b} f^2(x) dx - I^2 \right] \\
&= (b - a) \int_{a}^{b} M f(x) dx - (b - a) \int_{a}^{b} f^2(x) dx \\
&= (b - a) \int_{a}^{b} f(x) \big( M - f(x) \big) dx.
\end{align*}
Так как по условию $0 \le f(x) \le M$, подынтегральное выражение $f(x)(M - f(x)) \ge 0$ всюду на $[a,b]$, следовательно, интеграл неотрицателен:
$$\sigma^2_{geom} - \sigma^2_{mean} \ge 0 \implies \sigma^2_{mean} \le \sigma^2_{geom}.$$
Равенство достигается исключительно в тривиальном случае, когда $f(x) \in \{0, M\}$ почти всюду (кусочно-постоянная функция). Во всех практических расчетах геометрический метод заведомо проигрывает методу выборочного среднего по точности при равных затратах вычислений.
}

---

\subsubsection{Б. Методы понижения дисперсии (Variance Reduction)}

Погрешность ММК пропорциональна $\frac{\sigma}{\sqrt{N}}$. Экстенсивное увеличение $N$ вычислительно неэффективно ($N \sim \varepsilon^{-2}$). Единственный путь радикального повышения точности --- уменьшение теоретической дисперсии случайной величины $\sigma^2$.

\paragraph{1. Метод противоположных (антитетических) переменных (Antithetic Variates)}
Метод основан на использовании взаимно скоррелированных пар случайных величин с отрицательной ковариацией. Если случайная величина $U \sim \mathcal{U}(0,1)$, то величина $1 - U$ также строго распределена по закону $\mathcal{U}(0,1)$.

Вместо генерации $N$ независимых точек генерируются $N/2$ пар:
\begin{equation}
    X_i^{(1)} = a + (b - a) U_i, \quad X_i^{(2)} = a + (b - a) (1 - U_i), \quad i = 1, \dots, N/2.
\end{equation}

Элементарная оценка на одну пару:
\begin{equation}
    \theta_i = \frac{b - a}{2} \left[ f\big(X_i^{(1)}\big) + f\big(X_i^{(2)}\big) \right].
\end{equation}

Дисперсия полусуммы:
\begin{equation}
    \mathrm{Var}(\theta) = \frac{(b - a)^2}{4} \left[ \mathrm{Var}(f(X^{(1)})) + \mathrm{Var}(f(X^{(2)})) + 2\,\mathrm{Cov}\big(f(X^{(1)}), f(X^{(2)})\big) \right].
\end{equation}

Ввиду одинаковой распределенности $X^{(1)}$ и $X^{(2)}$:
\begin{equation}
    \mathrm{Var}(\theta) = \frac{(b - a)^2}{2} \mathrm{Var}(f(X)) + \frac{(b - a)^2}{2} \mathrm{Cov}\big(f(U), f(1 - U)\big).
\end{equation}

\infobox{Свойство отрицательной ковариации антитетических переменных}{
Если подынтегральная функция $f(x)$ монотонна (возрастает или убывает) на отрезке $[a,b]$, то функции $f(U)$ и $f(1 - U)$ имеют противоположные направления монотонности как функции переменной $U \in [0,1]$. Следовательно, их ковариация строго отрицательна:
$$\mathrm{Cov}\big(f(U), f(1 - U)\big) < 0.$$
Отрицательная ковариация частично компенсирует собственные дисперсии слагаемых, уменьшая результирующую дисперсию в десятки и сотни раз. Метод не требует дополнительных аналитических вычислений и тривиален в программной реализации.
}

\paragraph{2. Метод контрольных переменных (Control Variates)}
Пусть существует вспомогательная функция $h(x)$, качественно близкая к $f(x)$, для которой интеграл известен точно в аналитическом виде:
\begin{equation}
    I_h = \int_{a}^{b} h(x)\,dx.
\end{equation}

Интеграл $I$ тождественно представляется в виде:
\begin{equation}
    I = \int_{a}^{b} h(x)\,dx + \int_{a}^{b} \big[ f(x) - h(x) \big] dx = I_h + (b - a) \mathbb{E}\big[ f(X) - h(X) \big].
\end{equation}

Оценка ММК применяется исключительно к малой разности $\Delta(x) = f(x) - h(x)$:
\begin{equation}
    I_{CV} = I_h + \frac{b - a}{N} \sum_{i=1}^N \big[ f(X_i) - h(X_i) \big].
\end{equation}

Дисперсия остаточного члена:
\begin{equation}
    \mathrm{Var}(f(X) - h(X)) = \mathrm{Var}(f) + \mathrm{Var}(h) - 2\,\mathrm{Cov}(f, h).
\end{equation}
Если функция $h(x)$ аппроксимирует $f(x)$ с высоким коэффициентом корреляции ($\rho(f,h) \approx 1$), то дисперсия остатка стремится к нулю. В качестве $h(x)$ часто применяются отрезки ряда Тейлора или полиномы интерполяции малых степеней.

\paragraph{3. Выборка по важности (Importance Sampling)}
Пусть $g(x)$ --- вспомогательная плотность вероятности на $[a,b]$, удовлетворяющая условиям $g(x) > 0$ при $f(x) \ne 0$ и $\int_a^b g(x) dx = 1$. Перепишем интеграл:
\begin{equation}
    I = \int_{a}^{b} \frac{f(x)}{g(x)} g(x) dx = \mathbb{E}_g \left[ \frac{f(Y)}{g(Y)} \right], \quad Y \sim g(x).
\end{equation}

Оценка выборки по важности:
\begin{equation}
    I_{IS} = \frac{1}{N} \sum_{i=1}^N \frac{f(Y_i)}{g(Y_i)}, \quad Y_i \sim g(x).
\end{equation}

Дисперсия оценки:
\begin{equation}
    \mathrm{Var}\left( \frac{f(Y)}{g(Y)} \right) = \int_{a}^{b} \frac{f^2(x)}{g(x)} dx - I^2.
\end{equation}
Дисперсия обращается в абсолютный нуль ($\sigma^2 = 0$), если плотность идеальна:
\begin{equation}
    g^*(x) = \frac{|f(x)|}{\int_a^b |f(t)| dt} = \frac{f(x)}{I} \quad (\text{для } f(x) \ge 0).
\end{equation}
Хотя нормировочная константа $I$ неизвестна, выбор плотности $g(x)$, повторяющей основные геометрические черты графика $f(x)$, снижает дисперсию на порядки.

\paragraph{4. Расслоенная (стратифицированная) выборка (Stratified Sampling)}
Отрезок $[a,b]$ делится на $K$ равных частей $[x_{k-1}, x_k]$, на каждой из которых независимо генерируется по $N_k = N/K$ точек. По закону полной дисперсии общая дисперсия равна сумме средней внутригрупповой и межгрупповой дисперсий:
\begin{equation}
    \mathrm{Var}_{total} = \mathbb{E}[\mathrm{Var}(f|k)] + \mathrm{Var}(\mathbb{E}[f|k]).
\end{equation}
Стратификация полностью исключает межгрупповую вариацию $\mathrm{Var}(\mathbb{E}[f|k])$, ликвидируя эффект неравномерного скопления точек по отрезку.

\newpage

\section{Аналитический расчет практических заданий лабораторной работы №6}

\subsection{Задание 1.1: Вычисление интеграла $I = \int_0^1 e^x dx$}

\subsubsection{1. Точное аналитическое решение}
\begin{equation}
    I = \int_{0}^{1} e^x\,dx = \Big[ e^x \Big]_0^1 = e^1 - e^0 = e - 1 \approx 1.718281828459045.
\end{equation}

\subsubsection{2. Метод выборочного среднего}
Случайная величина $X \sim \mathcal{U}(0,1)$, $f(X) = e^X$.
\begin{align}
    \mathbb{E}[f(X)] &= \int_0^1 e^x dx = e - 1 \approx 1.718282, \\
    \mathbb{E}[f^2(X)] &= \int_0^1 e^{2x} dx = \left[ \frac{e^{2x}}{2} \right]_0^1 = \frac{e^2 - 1}{2} \approx 3.194528.
\end{align}
Теоретическая генеральная дисперсия:
\begin{equation}
    \sigma^2_{mean} = \mathbb{E}[f^2] - (\mathbb{E}[f])^2 = \frac{e^2 - 1}{2} - (e - 1)^2 = \frac{-e^2 + 4e - 3}{2} \approx 0.242036.
\end{equation}
Стандартное отклонение: $\sigma_{mean} = \sqrt{0.242036} \approx 0.491971$.

\subsubsection{3. Геометрический метод («Hit-or-Miss»)}
Функция монотонно возрастает на $[0,1]$, поэтому $M = \max_{[0,1]} e^x = e \approx 2.718282$.
Площадь прямоугольника $S_D = M \cdot 1 = e$.
Вероятность попадания: $p = \frac{I}{S_D} = \frac{e - 1}{e} \approx 0.632121$.
Теоретическая дисперсия геометрического метода:
\begin{equation}
    \sigma^2_{geom} = I(S_D - I) = (e - 1)\big(e - (e - 1)\big) = e - 1 \approx 1.718282.
\end{equation}
Соотношение дисперсий:
\begin{equation}
    \frac{\sigma^2_{geom}}{\sigma^2_{mean}} = \frac{1.718282}{0.242036} \approx 7.10.
\end{equation}
Геометрический метод дает дисперсию в 7.1 раза выше, чем метод выборочного среднего!

\subsubsection{4. Метод антитетических переменных}
Пара точек: $X^{(1)} = U$, $X^{(2)} = 1 - U$, где $U \sim \mathcal{U}(0,1)$.
Оценка пары: $\theta = \frac{e^U + e^{1 - U}}{2}$.
\begin{align}
    \mathbb{E}[\theta] &= e - 1, \\
    \mathbb{E}[\theta^2] &= \frac{1}{4} \mathbb{E}\big[ e^{2U} + 2e + e^{2(1 - U)} \big] = \frac{1}{4} \left( 2 \cdot \frac{e^2 - 1}{2} + 2e \right) = \frac{e^2 + 2e - 1}{4} \approx 2.956405.
\end{align}
Дисперсия антитетической оценки:
\begin{equation}
    \sigma^2_{anti} = \mathbb{E}[\theta^2] - (\mathbb{E}[\theta])^2 = 2.956405 - (1.718282)^2 \approx 0.003893.
\end{equation}
Коэффициент снижения дисперсии:
\begin{equation}
    \frac{\sigma^2_{mean}}{\sigma^2_{anti}} = \frac{0.242036}{0.003893} \approx 62.17.
\end{equation}
Дисперсия уменьшается более чем в 62 раза без усложнения вычислений!

\subsubsection{5. Метод контрольных переменных}
\begin{itemize}
    \item Линейная аппроксимация Тейлора $h_1(x) = 1 + x$: $I_{h_1} = \int_0^1 (1 + x)dx = 1.5$.
    Дисперсия остатка $\mathrm{Var}(e^X - (1+X)) \approx 0.04364$ (снижение в $5.55$ раз).
    \item Двухточечная интерполяция концов $h_2(x) = 1 + (e - 1)x$: $I_{h_2} = \frac{e + 1}{2} \approx 1.85914$.
    Дисперсия остатка $\mathrm{Var}(e^X - h_2(X)) \approx 0.00401$ (снижение в $60.3$ раз).
\end{itemize}

---

\subsection{Задание 1.2: Вычисление интеграла $I = \int_0^{\pi/2} \sin(x) dx$}

\subsubsection{1. Точное аналитическое решение}
\begin{equation}
    I = \int_{0}^{\pi/2} \sin(x)\,dx = \Big[ -\cos(x) \Big]_0^{\pi/2} = -\cos(\pi/2) + \cos(0) = 0 + 1 = 1.000000000000.
\end{equation}

\subsubsection{2. Метод выборочного среднего}
Отрезок $[0, \pi/2]$, длина $b - a = \pi/2$. Переменная $X = \frac{\pi}{2}U$, где $U \sim \mathcal{U}(0,1)$. Оценка $\xi = \frac{\pi}{2}\sin(X)$.
\begin{align}
    \mathbb{E}[\xi] &= 1.0, \\
    \mathbb{E}[\xi^2] &= \left(\frac{\pi}{2}\right)^2 \frac{2}{\pi} \int_0^{\pi/2} \sin^2(x) dx = \frac{\pi}{2} \cdot \frac{\pi}{4} = \frac{\pi^2}{8} \approx 1.233701.
\end{align}
Теоретическая генеральная дисперсия:
\begin{equation}
    \sigma^2_{mean} = \mathbb{E}[\xi^2] - (\mathbb{E}[\xi])^2 = \frac{\pi^2}{8} - 1.0 \approx 0.233701.
\end{equation}
Стандартное отклонение: $\sigma_{mean} = \sqrt{0.233701} \approx 0.483426$.

\subsubsection{3. Геометрический метод («Hit-or-Miss»)}
$M = \max \sin(x) = 1.0$. Площадь прямоугольника $S_D = 1 \cdot \frac{\pi}{2} = \frac{\pi}{2} \approx 1.570796$.
Вероятность попадания $p = \frac{1}{\pi/2} = \frac{2}{\pi} \approx 0.636620$.
Дисперсия геометрического метода:
\begin{equation}
    \sigma^2_{geom} = I(S_D - I) = 1 \cdot \left(\frac{\pi}{2} - 1\right) = \frac{\pi - 2}{2} \approx 0.570796.
\end{equation}
Соотношение дисперсий:
\begin{equation}
    \frac{\sigma^2_{geom}}{\sigma^2_{mean}} = \frac{0.570796}{0.233701} \approx 2.44.
\end{equation}

\subsubsection{4. Метод антитетических переменных}
Пара узлов: $X^{(1)} = \frac{\pi}{2}U$, $X^{(2)} = \frac{\pi}{2}(1 - U)$.
Оценка пары: $\theta = \frac{\pi}{4}\big[\sin(\frac{\pi}{2}U) + \sin(\frac{\pi}{2}(1 - U))\big]$.
Теоретическая дисперсия:
\begin{equation}
    \sigma^2_{anti} \approx 0.009563.
\end{equation}
Коэффициент снижения дисперсии:
\begin{equation}
    \frac{\sigma^2_{mean}}{\sigma^2_{anti}} = \frac{0.233701}{0.009563} \approx 24.44.
\end{equation}

\subsubsection{5. Метод контрольных переменных}
\begin{itemize}
    \item Линейная функция $h_1(x) = \frac{2}{\pi}x$: $I_{h_1} = \frac{\pi}{4} \approx 0.785398$. Дисперсия $\sigma^2_{CV_1} \approx 0.010131$ (снижение в $23.1$ раз).
    \item Параболическая аппроксимация $h_2(x) = \frac{4}{\pi}x - \frac{4}{\pi^2}x^2$:
    $$I_{h_2} = \int_0^{\pi/2} \left(\frac{4}{\pi}x - \frac{4}{\pi^2}x^2\right) dx = \frac{\pi}{3} \approx 1.047198.$$
    Дисперсия остатка:
    \begin{equation}
        \sigma^2_{CV_2} \approx 0.000940 \implies \frac{\sigma^2_{mean}}{\sigma^2_{CV_2}} \approx 248.7.
    \end{equation}
    Дисперсия падает почти в 250 раз!
\end{itemize}

\newpage

\section{Интегралы лабораторной работы №5 (Вариант 7) в контексте ММК}

В лабораторной работе №5 для \textbf{Варианта 7} были заданы следующие интегралы:

\subsection{Интеграл 1: $I_1 = \int_0^{\sqrt{3}} \sqrt{3 - x^2} dx$}

\subsubsection{1. Точное аналитическое решение}
Геометрически интеграл выражает площадь сектора круга радиуса $R = \sqrt{3}$ в первом квадранте (угол $\pi/2$):
\begin{equation}
    I_1 = \frac{1}{4} \pi R^2 = \frac{1}{4} \pi (\sqrt{3})^2 = \frac{3\pi}{4} \approx 2.356194490192345.
\end{equation}

\subsubsection{2. Метод выборочного среднего}
Отрезок $[0, \sqrt{3}]$, $b - a = \sqrt{3}$. Узел $X = \sqrt{3}U$, где $U \sim \mathcal{U}(0,1)$. Оценка $\xi = 3\sqrt{1 - U^2}$.
\begin{align}
    \mathbb{E}[\xi] &= 3 \int_0^1 \sqrt{1 - u^2} du = 3 \cdot \frac{\pi}{4} = \frac{3\pi}{4}, \\
    \mathbb{E}[\xi^2] &= 9 \int_0^1 (1 - u^2) du = 9 \left[ u - \frac{u^3}{3} \right]_0^1 = 9 \cdot \frac{2}{3} = 6.0.
\end{align}
Теоретическая дисперсия:
\begin{equation}
    \sigma^2_{mean} = 6.0 - \left(\frac{3\pi}{4}\right)^2 = 6 - \frac{9\pi^2}{16} \approx 0.448346.
\end{equation}
Стандартное отклонение: $\sigma_{mean} \approx 0.669586$.

\subsubsection{3. Антитетические переменные}
$\theta = \frac{3}{2}\big[\sqrt{1 - U^2} + \sqrt{1 - (1 - U)^2}\big]$.
Теоретическая дисперсия $\sigma^2_{anti} \approx 0.061813$. Снижение дисперсии в $7.25$ раза.

---

\subsection{Интеграл 2: $I_2 = \int_0^{3\pi/4} x \sin\left(\frac{2}{3}x\right) dx$}

\subsubsection{1. Точное аналитическое решение}
Интегрирование по частям: $u = x \implies du = dx$, $dv = \sin(2x/3)dx \implies v = -\frac{3}{2}\cos(2x/3)$.
\begin{align}
    I_2 &= \left[ -\frac{3}{2} x \cos\left(\frac{2}{3}x\right) \right]_0^{3\pi/4} + \frac{3}{2} \int_0^{3\pi/4} \cos\left(\frac{2}{3}x\right) dx \nonumber \\
    &= \left( -\frac{3}{2} \cdot \frac{3\pi}{4} \cdot \cos\left(\frac{\pi}{2}\right) - 0 \right) + \frac{3}{2} \left[ \frac{3}{2} \sin\left(\frac{2}{3}x\right) \right]_0^{3\pi/4} \nonumber \\
    &= 0 + \frac{9}{4} \sin\left(\frac{\pi}{2}\right) = \frac{9}{4} = 2.250000000000.
\end{align}

\subsubsection{2. Метод выборочного среднего}
Длина отрезка $B = \frac{3\pi}{4} \approx 2.356194$. Переменная $X = B \cdot U$. Оценка $\xi = B \cdot X \sin(2X/3)$.
Теоретическая дисперсия:
\begin{equation}
    \sigma^2_{mean} \approx 3.1987, \quad \sigma_{mean} \approx 1.7885.
\end{equation}

\subsubsection{3. Антитетические переменные}
Функция $f(x) = x \sin(2x/3)$ строго монотонно возрастает на $[0, 3\pi/4]$ (так как производная $f'(x) = \sin(2x/3) + \frac{2}{3}x\cos(2x/3) > 0$). Благодаря монотонности ковариация антитетических величин велика по модулю и отрицательна:
\begin{equation}
    \sigma^2_{anti} \approx 0.061409.
\end{equation}
Коэффициент снижения дисперсии:
\begin{equation}
    \frac{\sigma^2_{mean}}{\sigma^2_{anti}} = \frac{3.1987}{0.061409} \approx 52.09.
\end{equation}
Дисперсия падает более чем в 52 раза!

\newpage

\section{Сравнительный анализ: Лабораторная работа №6 vs Лабораторная работа №5}

\subsection{Сводная таблица сравнительных характеристик}

\begin{table}[h!]
\centering
\small
\begin{tabularx}{\textwidth}{p{3.8cm} X X}
\toprule
\textbf{Критерий сравнения} & \textbf{Лаб. работа №5 (Ньютон--Котес, Ричардсон)} & \textbf{Лаб. работа №6 (Метод Монте-Карло)} \\
\midrule
\textbf{Математическая природа} & Детерминированная кусочно-полиномиальная аппроксимация & Вероятностная точечная оценка математического ожидания \\
\addlinespace
\textbf{Характер результата} & Строго воспроизводимое детерминированное значение & Случайная величина, зависящая от псевдослучайного зерна (\texttt{seed}) \\
\addlinespace
\textbf{Сходимость в 1D ($d=1$)} & Трапеции: $O(N^{-2})$, \textbf{Симпсон: $O(N^{-4})$} & \textbf{$O(N^{-1/2})$} независимо от гладкости \\
\addlinespace
\textbf{Узлов для точности $10^{-4}$ в 1D} & Формула Симпсона: $N \sim 20$ узлов & ММК: $N \sim 10^7 - 10^8$ испытаний \\
\addlinespace
\textbf{Узлов для точности $10^{-8}$ в 1D} & Формула Симпсона: $N \sim 200$ узлов & ММК: $N \sim 10^{16}$ (практически недостижимо) \\
\addlinespace
\textbf{Сходимость в $d$-мерном кубе} & Трапеции: $O(N^{-2/d})$, Симпсон: $O(N^{-4/d})$ (\textit{проклятие размерности}) & \textbf{$O(N^{-1/2})$ всегда}, скорость не зависит от размерности $d$! \\
\addlinespace
\textbf{Порог эффективности по $d$} & Превосходит ММК при $d = 1, 2, 3$ & \textbf{Безусловно доминирует при $d \ge 4 \dots 8$} \\
\addlinespace
\textbf{Требования к гладкости $f(x)$} & Симпсон требует $f \in C^4$, трапеции $f \in C^2$ & Требуется лишь $f \in L_2$ (конечность $\int f^2 dx$) \\
\addlinespace
\textbf{Влияние краевых сингулярностей (Вариант 7)} & На $\sqrt{3-x^2}$ из-за $f' \to -\infty$ порядок Симпсона упал до $O(h^{1.5})$ & \textbf{Полный иммунитет:} сохраняет скорость $O(N^{-1/2})$ без аномалий \\
\addlinespace
\textbf{Сложность геометрии границы $\Omega$} & Требует параметризации и триангуляции & Тривиально: предикат \texttt{is\_inside(x)} \\
\addlinespace
\textbf{Оценка погрешности} & Правило Рунге (требует пересчета на сетке с шагом $h/2$) & «На лету» в один проход через выборочную дисперсию $s / \sqrt{N}$ \\
\addlinespace
\textbf{Гарантия погрешности} & Строгая детерминированная граница $|R_n| \le C h^p$ & Вероятностная (доверительный интервал с надежностью $1-\alpha$) \\
\addlinespace
\textbf{Параллелизуемость} & Требует согласования границ подынтервалов & Идеальная независимость потоков (\textit{embarrassingly parallel}) \\
\bottomrule
\end{tabularx}
\caption{Сравнительный анализ методов квадратур Ньютона--Котеса и метода Монте-Карло}
\end{table}

---

\subsection{1. Одномерное интегрирование: превосходство квадратур Ньютона--Котеса}
В одномерном случае ($d=1$) для гладких функций квадратурные формулы Ньютона--Котеса и Гаусса абсолютно непревзойденны:
\begin{itemize}
    \item Для интеграла $I_2 = \int_0^{3\pi/4} x \sin(2x/3) dx = 2.25$: формуле Симпсона достаточно $N = 32$ узла, чтобы дать погрешность менее $10^{-10}$ за доли микросекунды.
    \item Методу Монте-Карло для аналогичной точности потребовалось бы астрономическое число точек $N \sim 10^{20}$, что заняло бы годы непрерывных вычислений на современном суперкомпьютере.
    \item Способы понижения дисперсии (антитетические переменные, контрольные переменные) уменьшают постоянный множитель при $N^{-1/2}$ в 10--250 раз, но \textbf{не способны изменить фундаментальный степенной показатель $-1/2$}.
\end{itemize}

---

\subsection{2. Многомерное интегрирование: проклятие размерности и торжество ММК}
В $d$-мерном единичном кубе $[0,1]^d$:
\begin{itemize}
    \item Чтобы построить регулярную сетку с шагом $h = 1/m$, требуется суммарно $N = m^d$ узлов. Шаг выражается через число точек как $h = N^{-1/d}$.
    \item Погрешность составной формулы Симпсона равна $O(h^4) = O(N^{-4/d})$.
    \item Погрешность составной формулы трапеций равна $O(h^2) = O(N^{-2/d})$.
    \item Погрешность метода Монте-Карло равна $O(N^{-1/2})$ независимо от $d$.
\end{itemize}

Сравнение показателей степени:
\begin{itemize}
    \item При $d = 4$: формула трапеций дает $N^{-2/4} = N^{-1/2}$ --- точность трапеций сравнивается с базовым Монте-Карло.
    \item При $d = 8$: формула Симпсона дает $N^{-4/8} = N^{-1/2}$ --- Симпсон сравнивается с Монте-Карло.
    \item При $d > 8$: метод Симпсона сходится \textbf{медленнее}, чем простейший метод Монте-Карло!
    \item При $d = 20$: для простейшей сетки всего по $m = 10$ точек по каждому направлению требуется $N = 10^{20}$ точек (невозможно сохранить в памяти и обсчитать за разумное время). Метод Монте-Карло позволяет найти интеграл с точностью $1\%$ за $10^6$ вызовов функции за пару секунд.
\end{itemize}

---

\subsection{3. Влияние гладкости подынтегральной функции: фундаментальный урок Варианта 7}
В лабораторной работе №5 на Интеграле 1 Варианта 7:
\begin{equation}
    I_1 = \int_0^{\sqrt{3}} \sqrt{3 - x^2} dx
\end{equation}
проявилась фундаментальная уязвимость полиномиальных квадратур:
\begin{itemize}
    \item В точке $x \to \sqrt{3}$ первая производная обращается в бесконечность: $f'(x) = -\frac{x}{\sqrt{3-x^2}} \to -\infty$.
    \item Теоретическая оценка погрешности формулы Симпсона базируется на ограничении четвертой производной: $|R_n| \le \frac{(b-a)^5}{2880} \max |f^{(4)}(\xi)| h^4$. Из-за сингулярности производной реальный порядок сходимости метода Симпсона в лабораторной работе №5 снизился с теоретического $O(h^4)$ до фактического $O(h^{1.5})$, а экстраполяция Ричардсона давала сбои.
    \item \textbf{Для метода Монте-Карло производные функции не имеют значения.} Единственное условие применимости ЦПТ --- существование конечной генеральной дисперсии, то есть интегрируемость квадрата функции:
    $$\int_0^{\sqrt{3}} \big(\sqrt{3 - x^2}\big)^2 dx = \int_0^{\sqrt{3}} (3 - x^2) dx = 2\sqrt{3} < \infty.$$
    Поскольку дисперсия конечна ($\sigma^2_{mean} \approx 0.448$), метод Монте-Карло демонстрирует безукоризненную сходимость $O(N^{-1/2})$ без каких-либо отклонений.
\end{itemize}

---

\subsection{4. Способ оценки и гарантия погрешности}
\begin{itemize}
    \item \textbf{Лабораторная работа №5:} Оценка погрешности производилась апостериорно по правилу Рунге:
    $$\varepsilon_{Runge} = \frac{I(h/2) - I(h)}{2^p - 1}.$$
    Это требовало проведения полного повторного расчета на сгущенной сетке с шагом $h/2$, что увеличивало объем вычислений вдвое. Оценка носила детерминированный характер.
    \item \textbf{Лабораторная работа №6:} Оценка погрешности формируется параллельно с накоплением суммы интеграла через выборочную дисперсию $s^2$. Дополнительных вычислений функции не требуется. При этом оценка является стохастической (гарантия с вероятностью $95\%$ или $99.7\%$).
\end{itemize}

\newpage

\section{Рекомендации к программной реализации на C++}

\subsection{1. Корректный генератор псевдослучайных чисел}
В современном стандарте C++11/17/20 рекомендуется использовать 64-битный вихрь Мерсенна:
\begin{lstlisting}
#include <random>
#include <cmath>

std::random_device rd;
std::mt19937_64 rng(rd()); // 64-bit Mersenne Twister engine
std::uniform_real_distribution<double> dist01(0.0, 1.0);
\end{lstlisting}

\subsection{2. Численно устойчивое накопление статистики (Алгоритм Велфорда)}
Прямое вычисление выборочной дисперсии по стандартной формуле $\sum x_i^2 - \frac{1}{N}(\sum x_i)^2$ при больших $N \ge 10^7$ приводит к катастрофической потере точности из-за вычитания близких чисел. Рекомендуется использовать однопроходный численно стабильный алгоритм Велфорда:

\begin{lstlisting}
struct MonteCarloAccumulator {
    uint64_t count = 0;
    double mean = 0.0;
    double M2 = 0.0; // Sum of squared deviations from the current mean

    void update(double x) {
        count++;
        double delta = x - mean;
        mean += delta / count;
        double delta2 = x - mean;
        M2 += delta * delta2;
    }

    double get_mean() const { return mean; }
    double get_variance() const { return count > 1 ? M2 / (count - 1) : 0.0; }
    double get_std_error() const { return std::sqrt(get_variance() / count); }
    double get_ci95() const { return 1.95996 * get_std_error(); }
    double get_ci997() const { return 3.00000 * get_std_error(); }
};
\end{lstlisting}

\subsection{3. Пример реализации методов понижения дисперсии}
\begin{lstlisting}
// 1. Integral exp(x) on [0, 1], Sample Mean method
double u = dist01(rng);
double val_mean = std::exp(u);

// 2. Integral exp(x) on [0, 1], Antithetic Variates
double u = dist01(rng);
double val_anti = 0.5 * (std::exp(u) + std::exp(1.0 - u));

// 3. Integral exp(x) on [0, 1], Control Variate h(x) = 1 + (e - 1)*x
double u = dist01(rng);
double h = 1.0 + (std::numbers::e - 1.0) * u;
double val_cv = (std::numbers::e + 1.0) / 2.0 + (std::exp(u) - h);
\end{lstlisting}

\section{Заключение}
Метод Монте-Карло представляет собой универсальный численный аппарат, обладающий принципиальными достоинствами при решении многомерных задач ($d \ge 4$), задач со сложной геометрией и задач с недифференцируемыми подынтегральными функциями. В одномерных задачах для гладких функций он значительно уступает по скорости квадратурам Ньютона--Котеса (Лабораторная работа №5). Однако методы понижения дисперсии (антитетические и контрольные переменные) позволяют снизить дисперсию в десятки и сотни раз, существенно расширяя область эффективного применения стохастических методов в вычислительной физике.

\end{document}
